% !PRINTED 31
\weberpage{31}

Demnach haben wir auch die Formel
\begin{equation}
 {f(x)-f(\alpha)\over x-\alpha}=Q(x),
\tag{7}
\end{equation}
worin $Q(x)$ eine ganze Function vom Grade $n-1$ ist.

Wenn wir in den Ausdrücken (6) an Stelle der unbestimmten Grösse
$\alpha$ das Zeichen $x$ setzen, so entsteht daraus eine Reihe von
ganzen rationalen Functionen von $x$, die wir, wenn wir der Einfachheit
halber $a_0=1$ setzen, so schreiben:
\begin{equation}
\begin{aligned}
 f_0&=1,\\
 f_1&=x+a_1,\\
 f_2&=x^2+a_1x+a_2,\\
 &\cdots,\\
 f_{n-1}&=x^{n-1}+a_1x^{n-2}+a_2x^{n-3}+\cdots+a_{n-1}.
\end{aligned}
\tag{8}
\end{equation}
Diese Functionen $f_0,f_1,\ldots,f_{n-1}$ werden uns später noch gute
Dienste leisten. Für jetzt fügen wir noch folgende Bemerkungen bei.

Man kann nach (8) die Potenzen $1,x,x^2,\ldots,x^{n-1}$ von $x$
linear ausdrücken durch die Functionen $f_0,f_1,\ldots,f_{n-1}$, und
zwar so, dass in den Coëfficienten nur ganze rationale Verbindungen der
$a$ vorkommen, z. B.
\[
\begin{aligned}
 1&=f_0,\\
 x&=f_1-a_1f_0,\\
 x^2&=f_2-a_1f_1+(a_1^2-a_2)f_0,\\
 &\cdots,
\end{aligned}
\]
und daraus folgt, dass man jede ganze rationale Function von $x$,
deren Grad nicht grösser als $n-1$ ist, gleichfalls linear durch
$f_0,f_1,\ldots,f_{n-1}$ ausdrücken kann in der Form
\begin{equation}
 y_0f_0+y_1f_1+\cdots+y_{n-1}f_{n-1},
\tag{9}
\end{equation}
worin die Coëfficienten $y_0,y_1,\ldots,y_{n-1}$ von $x$ unabhängig
sind.

Ist also $F(x)$ eine beliebige ganze rationale Function von $x$, so
kann man nach \S~3, indem man $f(x)$ als Divisor betrachtet,
\begin{equation}
 F(x)=Qf(x)+y_0f_0+y_1f_1+\cdots+y_{n-1}f_{n-1}
\tag{10}
\end{equation}
setzen, worin auch $Q$ eine ganze rationale Function von $x$ ist.

Zur recurrenten Berechnung der Functionen $f_r(x)$ ergiebt sich aus
(8) die Relation:
\begin{equation}
 f_r(x)-xf_{r-1}(x)=a_r.
\tag{11}
\end{equation}
